\section{R ile çözümlü uygulama}
\label{sec:r-anova}

\textbf{Soru.} Bölümdeki üç grubun $n=(10,10,10)$,
$\bar x=(70,75,82)$ ve $s=(6,7,5)$ özetlerinden klasik ANOVA'yı,
etki büyüklüklerini ve Tukey eşzamanlı fark aralıklarını hesaplayın.
Özetleri ham veri gibi \texttt{aov} komutuna vermeyin.

\begin{lstlisting}[style=prismR]
hacim <- c(10, 10, 10)
ortalama <- c(70, 75, 82)
standart_sapma <- c(6, 7, 5)
grup_sayisi <- length(hacim)
genel <- weighted.mean(ortalama, hacim)
ss_grup <- sum(hacim * (ortalama - genel)^2)
ss_hata <- sum((hacim - 1) * standart_sapma^2)
sd_grup <- grup_sayisi - 1
sd_hata <- sum(hacim) - grup_sayisi
mse <- ss_hata / sd_hata
f_degeri <- (ss_grup / sd_grup) / mse
c(F = f_degeri,
  p = pf(f_degeri, sd_grup, sd_hata, lower.tail = FALSE),
  eta2 = ss_grup / (ss_grup + ss_hata),
  omega2 = (ss_grup - sd_grup * mse) /
    (ss_grup + ss_hata + mse))
ciftler <- combn(seq_along(hacim), 2)
kritik <- qtukey(0.95, nmeans = grup_sayisi, df = sd_hata)
tukey <- t(apply(ciftler, 2, function(cift) {
  fark <- ortalama[cift[1]] - ortalama[cift[2]]
  se <- sqrt(mse / 2 * sum(1 / hacim[cift]))
  c(fark = fark, alt = fark - kritik * se,
    ust = fark + kritik * se,
    p_duzeltilmis = ptukey(abs(fark) / se,
      nmeans = grup_sayisi, df = sd_hata, lower.tail = FALSE))
}))
rownames(tukey) <- c("A-B", "A-C", "B-C")
tukey
\end{lstlisting}

\begin{outputguidebox}
\textbf{Çözüm ve kontrol değerleri.} Genel ortalama 75,6667;
$SS_{\text{grup}}=726{,}6667$, $SS_{\text{hata}}=990$ ve
$MSE=36{,}6667$'dir. $F(2,27)=9{,}9091$, $p=0{,}0005927$,
$\eta^2=0{,}4233$, $\omega^2=0{,}3726$ bulunur. Tukey aralıklarının
yarı genişliği 6,7143'tür. A--B aralığı $[-11{,}7143;1{,}7143]$,
A--C aralığı $[-18{,}7143;-5{,}2857]$, B--C aralığı
$[-13{,}7143;-0{,}2857]$'dir. Son iki aralık sıfırı dışlar.
\end{outputguidebox}

\textbf{Yorum.} Klasik ANOVA ve bu Tukey karşılaştırmaları bağımsız
gözlemler, grup içi yaklaşık normallik ve ortak hata varyansı modeline dayanır.
Genel testin anlamlılığı, üç çiftin hepsinin farklı olduğunu söylemez.
Ham veri bulunduğunda \texttt{aov(puan \textasciitilde{} grup, data = veri)}
ve ardından \texttt{TukeyHSD} kullanılabilir; grup değişkeni faktör olmalıdır.

\subsection*{Eşit olmayan varyanslarda ek çözüm}
Bölümdeki Welch örneğinin özetleriyle aynı düzeltmeyi doğrudan hesaplayın.

\begin{lstlisting}[style=prismR]
hacim <- c(12, 20, 8)
ortalama <- c(70, 76, 84)
varyans <- c(3, 10, 5)^2
grup_sayisi <- length(hacim)
agirlik <- hacim / varyans
pay <- agirlik / sum(agirlik)
merkez <- sum(pay * ortalama)
duzeltme <- sum((1 - pay)^2 / (hacim - 1))
f_welch <- (sum(agirlik * (ortalama - merkez)^2) /
  (grup_sayisi - 1)) /
  (1 + 2 * (grup_sayisi - 2) * duzeltme /
     (grup_sayisi^2 - 1))
sd_payda <- (grup_sayisi^2 - 1) / (3 * duzeltme)
c(F = f_welch, sd_pay = grup_sayisi - 1,
  sd_payda = sd_payda,
  p = pf(f_welch, grup_sayisi - 1, sd_payda,
         lower.tail = FALSE))
\end{lstlisting}

\textbf{Çözüm.} Yaklaşık $F=25{,}3265$, pay serbestlik derecesi 2 ve
payda serbestlik derecesi 18,1621 bulunur; $p=0{,}000005578$'dir. Ham veride
\texttt{oneway.test(puan \textasciitilde{} grup, data = veri, var.equal = FALSE)}
aynı yöntemi uygular. Klasik ANOVA'nın ortak MSE'siyle hesaplanan Tukey
aralıklarını bu eşit olmayan varyans örneğine taşımayın.

\textbf{Pekiştirme ve yanıt.} Klasik örnekte A--B farkı $-5$ olsa da
aralık sıfırı içerir. Bu, iki yöntemin eşdeğer olduğunun kanıtı değildir;
yüzde 5 aile düzeyinde fark gösterilemediğini söyler.